\chapter{Begleitende Materialien}
\label{app:begleitmaterial}

Dieser Anhang dokumentiert den Inhalt des dieser Arbeit beigelegten ZIP-Archivs.
Das Archiv enthält den im Rahmen dieser Arbeit entwickelten Quellcode, die zugehörige Dokumentation sowie die aggregierte Resultattabelle \texttt{ExperimentLogs.csv}.
Es ermöglicht, die in Kapitel~\ref{cha:realisierung} beschriebene Pipeline nachzuvollziehen und die in Kapitel~\ref{cha:diskussion} dargestellten Resultate auf einer vergleichbaren Maschine zu reproduzieren.
Bewusst nicht enthalten sind grosse oder von Dritten stammende Bestandteile wie rohe Verkehrsmitschnitte, Datensätze, trainierte Modelle und die WFlib-Codebasis, die sich aus dem enthaltenen Quellcode reproduzieren lassen und in Abschnitt~\ref{app:begleitmaterial_ausschluesse} einzeln aufgeführt sind.
Die finale Bachelorarbeit als PDF und ihre LaTeX-Quellen werden separat abgegeben und sind nicht Teil dieses Code-Archivs.

\section{Verzeichnisstruktur}
\label{app:begleitmaterial_struktur}

Tabelle~\ref{tab:begleitmaterial_struktur} listet die wichtigsten Pfade des Archivs und deren Inhalt.
Die Pfadangaben sind relativ zur Wurzel des Archivs und entsprechen der Struktur des Projekt-Repositoriums, sodass die relativen Pfade in den Skripten und in \texttt{README.md} unverändert gelten.

\begin{table}[H]
\centering
\caption{Inhalt des beigelegten ZIP-Archivs.}
\label{tab:begleitmaterial_struktur}
\begin{tabular}{p{5.0cm}p{8.5cm}}
\hline
Pfad & Inhalt \\
\hline
\texttt{README.md} & Einstiegspunkt mit Aufbau, Bedienung und Reproduktionsanleitung \\
\texttt{documentation.md} & Vertiefende Gesamtbeschreibung der Pipeline (Setup, Ablauf, Architektur) \\
\texttt{ExperimentLogs.csv} & Aggregierte Resultattabelle aller Experimente \\
\texttt{src/simulation/} & Skripte zur Steuerung der Shadow-Simulation (Abschnitt~\ref{sec:realisierung_uebersicht}) \\
\texttt{src/simulation/conf/} & torrc-Vorlagen und Tamaraw-Bridge-Zertifikat \\
\texttt{src/simulation/generated/urls.txt} & Aus der \gls{zim}-Datei erzeugte URL-Liste der Zielseiten \\
\texttt{src/ml/} & Skripte der ML-Pipeline (\texttt{pcap\_to\_npz.py}, \texttt{run\_df.sh}, \texttt{correlation.py}, \texttt{pcap\_overhead.py}, \texttt{plot\_*.py}, \texttt{report.py}) \\
\hline
\end{tabular}
\end{table}

\section{Nicht enthaltene Daten}
\label{app:begleitmaterial_ausschluesse}

Aus Platzgründen und zur Wahrung von Lizenz- und Urheberrechten werden die in Tabelle~\ref{tab:begleitmaterial_ausschluesse} aufgeführten Bestandteile nicht im Archiv ausgeliefert.
Für jeden Eintrag wird der Grund und der reproduzierende Schritt angegeben.

\begin{table}[H]
\centering
\caption{Bewusst nicht ausgelieferte Bestandteile und ihre Reproduktion.}
\label{tab:begleitmaterial_ausschluesse}
\begin{tabular}{p{3.7cm}p{4.2cm}p{5.4cm}}
\hline
Bestandteil & Grund & Reproduktion \\
\hline
Rohe \glspl{pcap} der Monitor- und Exit-Hosts & Insgesamt rund 250 GB, sprengt die übliche Archivgrösse & \texttt{python3 shadowctl.py run <name>} mit der jeweiligen Konfiguration (Abschnitt~\ref{app:begleitmaterial_reproduktion}) \\
Experiment-Konfigurationen (\texttt{exp-*/} mit \texttt{shadow.config.yaml}, \texttt{schedule.json}, \texttt{sim.log}) & Werden bei jedem Lauf aus den enthaltenen Skripten neu erzeugt & Entstehen automatisch bei \texttt{shadowctl.py run} \\
Datensätze (\texttt{.npz}), trainierte \gls{df}-Modelle und Ergebnis-JSONs & Aus den \glspl{pcap} ableitbar, keine Primärdaten & \texttt{pcap\_to\_npz.py} und \texttt{run\_df.sh} (Abschnitt~\ref{app:begleitmaterial_reproduktion}) \\
Python-Virtualenvs (\texttt{venv/}) & Maschinen- und versionsspezifisch & Neu anlegen via \texttt{python3 -m venv venv} und Pakete aus \texttt{documentation.md} \\
WFlib-Codebasis (\texttt{src/ml/wflib/}) & Drittprojekt unter eigener Lizenz~\cite{deng2023robust} & Klonen aus \texttt{https://github.com/Xinhao-Deng/Website-Fingerprinting-Library} \\
Tornettools-Eingabedaten (Konsensus, Server-Descriptors) & Drittprojekt, gross und versionsspezifisch & Bezug über \texttt{https://github.com/shadow/tornettools} und das dort dokumentierte CollecTor-Snapshot-Verfahren \\
\hline
\end{tabular}
\end{table}

\section{Reproduktion eines Resultats}
\label{app:begleitmaterial_reproduktion}

Zur Veranschaulichung wird hier exemplarisch die Reproduktion des \gls{closedworld}-Laufs \texttt{exp-baseline-80} skizziert, dessen Ergebnis in Tabelle~\ref{tab:experimentlogs} dokumentiert ist.

\begin{enumerate}
    \item Simulation starten und Pcaps erzeugen: aus \texttt{src/simulation/} den Befehl
        \texttt{python3 shadowctl.py run exp-baseline-80 --pages 80 --visits 150 --monitors 20 --padding off}
        absetzen und die Pcaps anschliessend mit \texttt{python3 shadowctl.py pull-results --name exp-baseline-80} herunterladen.
    \item Pcaps in das WFlib-Format konvertieren: aus \texttt{src/ml/} den Befehl
        \texttt{python3 pcap\_to\_npz.py --schedule ../simulation/exp-baseline-80/shadow.config.schedule.json --shadow-data ../simulation/exp-baseline-80/results/shadow.data --out wflib/datasets/ExpBaseline80.npz}
        ausführen.
    \item Daten splitten, trainieren und testen: aus \texttt{src/ml/} das Wrapper-Skript \texttt{bash run\_df.sh ExpBaseline80} aufrufen, das die in \texttt{wflib/exp/} liegenden WFlib-Skripte mit den in dieser Arbeit verwendeten Hyperparametern startet.
    \item Resultate ablesen: \texttt{cat wflib/logs/ExpBaseline80/DF/result.json} liefert die in Tabelle~\ref{tab:experimentlogs} verglichenen Kennzahlen.
\end{enumerate}

Da das Archiv aus Platzgründen weder vorkonvertierte Datensätze noch trainierte Modelle enthält, durchläuft die Reproduktion die vollständige Kette von der Simulation über die Konvertierung bis zum Testlauf.
Die dafür benötigte WFlib-Codebasis und das Python-Environment werden zuvor gemäss \texttt{README.md} und \texttt{documentation.md} eingerichtet.

\section{Tabellarische Gesamtsicht aller Hauptexperimente}
\label{app:experimentlogs}

Tabelle~\ref{tab:experimentlogs} fasst die in Kapitel~\ref{cha:diskussion} ausgewerteten zwölf Hauptexperimente in einer einzigen Übersicht zusammen.
Die Werte entsprechen direkt der Datei \texttt{ExperimentLogs.csv} im Hauptverzeichnis des begleitenden Repositoriums und sind auf die in dieser Arbeit ausgewerteten Spalten reduziert.
Konstante Spalten wie Netzwerk-Skala, Seed, \gls{zim}-Datei oder Trainings-Epochen werden weggelassen, da sie für alle Hauptexperimente identisch sind (Network Scale 0.010, Seed 1, \gls{zim} \texttt{simple\_en\_maxi\_2026-02}, 30 Trainings-Epochen, Visit-Interval 30~s, Visits pro Seite 150).
Explorative Läufe (Smoketests, frühe Korrelations- und \gls{tamaraw}-Vorversuche) sind nicht enthalten, sind in der vollständigen CSV-Datei aber dokumentiert.

\begin{landscape}
\footnotesize
\setlength{\tabcolsep}{4pt}
\begin{longtable}{lrlrrrrrrrrrr}
\caption{Aggregierte Kennzahlen der zwölf Hauptexperimente. Quelle: \texttt{ExperimentLogs.csv}.} \label{tab:experimentlogs} \\
\hline
Experiment & Klassen & Padding & Wall-Clock & Samples & /Klasse & Fail\,\% & Pkt/Sample & Acc.\,\% & Prec.\,\% & Rec.\,\% & F1\,\% & Rand.\,\% \\
\hline
\endfirsthead
\multicolumn{13}{l}{\textit{Fortsetzung von Tabelle~\ref{tab:experimentlogs}}} \\
\hline
Experiment & Klassen & Padding & Wall-Clock & Samples & /Klasse & Fail\,\% & Pkt/Sample & Acc.\,\% & Prec.\,\% & Rec.\,\% & F1\,\% & Rand.\,\% \\
\hline
\endhead
\hline
\multicolumn{13}{r}{\textit{Fortsetzung nächste Seite}} \\
\endfoot
\hline
\endlastfoot
exp-baseline-20       &  20 & OFF     & 07:57 &  3\,000 & 150 & 1.7 &    635 & 92.7 & 93.3 & 92.7 & 92.7 & 5.0 \\
exp-padding-20        &  20 & ON      & 08:16 &  3\,000 & 150 & 1.5 &    638 & 83.3 & 84.2 & 83.3 & 82.8 & 5.0 \\
exp-reduced-20        &  20 & Reduced & 08:18 &  3\,000 & 150 & 1.9 &    611 & 86.7 & 88.1 & 86.7 & 86.3 & 5.0 \\
exp-tamaraw-20v2      &  20 & Tamaraw & 08:19 &  3\,000 & 150 & 1.6 & 5\,000 & 39.3 & 50.7 & 39.3 & 38.6 & 5.0 \\
exp-baseline-80       &  80 & OFF     & 15:48 & 12\,000 & 150 & 4.4 &    583 & 82.4 & 81.8 & 82.4 & 81.4 & 1.2 \\
exp-padding-80        &  80 & ON      & 16:01 & 12\,000 & 150 & 4.5 &    601 & 78.6 & 78.5 & 78.6 & 76.9 & 1.2 \\
exp-reduced-80        &  80 & Reduced & 15:39 & 12\,000 & 150 & 4.1 &    591 & 77.5 & 76.7 & 77.5 & 76.2 & 1.2 \\
exp-tamaraw-80        &  80 & Tamaraw & 16:47 & 12\,000 & 150 & 2.9 & 5\,000 & 31.6 & 39.4 & 31.6 & 29.2 & 1.2 \\
exp-openworld         & 280 & OFF     & 17:48 & 13\,740 & 146 & 3.1 &    601 & 68.3 & 67.9 & 74.6 & 69.8 & 1.2 \\
exp-openworld-on      & 280 & ON      & 18:00 & 13\,680 & 146 & 3.3 &    604 & 70.7 & 72.5 & 75.6 & 72.5 & 1.2 \\
exp-openworld-reduced & 280 & Reduced & 18:03 & 13\,680 & 146 & 3.1 &    592 & 71.9 & 73.3 & 76.4 & 73.6 & 1.2 \\
exp-tamaraw-openworld & 280 & Tamaraw & 16:45 & 12\,620 & 132 & 3.2 & 4\,994 & 28.0 & 26.6 & 21.4 & 20.0 & 1.2 \\
\end{longtable}
\normalsize
\end{landscape}

Die Spalten haben folgende Bedeutung.
Klassen entspricht der Anzahl der vom Klassifikator unterschiedenen Seiten (\gls{closedworld}) beziehungsweise der Anzahl Quellseiten (\gls{openworld}: 280 Quellseiten, davon 80 überwachte Klassen plus eine Sammelklasse).
Padding bezeichnet die im Tor-Daemon der \gls{monitorhost}s gesetzte Verteidigung (OFF, ON, Reduced) beziehungsweise das auf den Bridge-Hosts aktivierte \gls{tamaraw}.
Wall-Clock ist die Laufzeit der \gls{shadow}-Simulation in Stunden:Minuten.
Samples und /Klasse geben die nach Konvertierung mit \texttt{pcap\_to\_npz.py} verfügbaren Aufzeichnungen wieder.
Fail\,\% beziffert den Anteil der durch \texttt{wget2} abgebrochenen Seitenabrufe.
Pkt/Sample ist die mittlere Anzahl Pakete pro Aufzeichnung über die ersten 5000 Werte des direction-Vektors.
Acc., Prec., Rec.\ und F1 stammen aus dem WFlib-Testlauf des \gls{df}-Modells.
Rand.\ ist die theoretische Zufallsbasis $1/K$ mit $K$ effektiv vom Klassifikator unterschiedenen Klassen (\gls{openworld}: 81 = 80 überwachte + 1 Sammelklasse, daher 1.2\,\%).

\section{Verteilung der Sequenzlänge pro Besuch}
\label{app:seqlen}

Tabelle~\ref{tab:eval_seqlen} ergänzt die in Abschnitt~\ref{sec:eval_validierung} herangezogenen Sanity-Kontrollen um die empirische Verteilung der Sequenzlänge pro Besuch und dokumentiert, in welchem Umfang die in Abschnitt~\ref{sec:konzepte_features} festgelegte Obergrenze von 5000~Werten greift.

\begin{table}[H]
    \centering
    \caption{Verteilung der Sequenzlänge pro Besuch nach der Konvertierung mit \texttt{pcap\_to\_npz.py}. Eigene Auswertung über alle Samples des jeweiligen Datensatzes.}
    \label{tab:eval_seqlen}
    \begin{tabular}{lrrrrr}
        \hline
        Experiment & N & Median & p90 & Max & Saturiert ($\geq$5000) \\
        \hline
        exp-baseline-20       &  3\,000 &    623 &    854 & 1\,270 &   0.0\,\% \\
        exp-padding-20        &  3\,000 &    623 &    832 & 1\,247 &   0.0\,\% \\
        exp-reduced-20        &  3\,000 &    596 &    827 & 1\,224 &   0.0\,\% \\
        exp-tamaraw-20v2      &  3\,000 & 5\,000 & 5\,000 & 5\,000 & 100.0\,\% \\
        exp-baseline-80       & 12\,000 &    581 &    733 & 2\,303 &   0.0\,\% \\
        exp-padding-80        & 12\,000 &    601 &    734 & 2\,253 &   0.0\,\% \\
        exp-reduced-80        & 12\,000 &    591 &    736 & 2\,311 &   0.0\,\% \\
        exp-tamaraw-80        & 12\,000 & 5\,000 & 5\,000 & 5\,000 & 100.0\,\% \\
        exp-openworld         & 13\,740 &    597 &    747 & 2\,399 &   0.0\,\% \\
        exp-openworld-on      & 13\,680 &    599 &    732 & 2\,270 &   0.0\,\% \\
        exp-openworld-reduced & 13\,680 &    588 &    734 & 2\,454 &   0.0\,\% \\
        exp-tamaraw-openworld & 12\,620 & 5\,000 & 5\,000 & 5\,000 &  99.8\,\% \\
        \hline
    \end{tabular}
\end{table}
